% =====================================================================
\chapter[From Intent to Verifiable Execution]{From Intent to Verifiable Execution}\label{ch:intent}
% =====================================================================

This chapter addresses Problem \ref{prob:trust}: if the cost of coding tends to zero and no
human reviews the implementation line by line, how is correctness guaranteed
mechanically. The answer has three parts: decompose correctness into dimensions that can be
treated separately, define the intermediate process as a verifiable refinement chain, and
organise trust as a composite of two modes of generation. The chapter also identifies an
entry point on that chain which cannot be mechanically closed, and the sense in which
consensus can approximate it.

\section{Dimensions of Correctness}\label{sec:five-correctness}

An AI rewrites a 1\,ns call into a 1\,$\mu$s call: functionally identical, but a thousand
times slower. The example shows that ``correct'' is not a single dimension.

\begin{definition}[The five dimensions of correctness]\label{def:five-correctness}
See Table \ref{tab:five-correctness}.
\end{definition}

\begin{table}[htbp]
\centering
\small
\caption{The five dimensions of correctness and their mechanical checkability}\label{tab:five-correctness}
\begin{tabular}{L{2.3cm}L{6.6cm}L{5.5cm}}
\toprule
\textbf{Dimension} & \textbf{Question to be answered} & \textbf{Mechanically checkable?} \\
\midrule
Functional correctness & Do the outputs and the state changes meet the requirements? &
  For an explicit specification, provable or testable \\
Side-effect correctness & Has the program performed unauthorised I/O, writes, network
  requests or state changes? &
  Constrainable by effect specifications, permission boundaries and runtime isolation \\
Resource correctness & Do time, memory, bandwidth and energy meet the budget? &
  Partly provable statically; partly measurable only, or statistically guaranteed in a
  specified environment \\
System correctness & After components are composed, are concurrency, failure recovery
  and state consistency still correct? &
  Partly formally provable; partly requiring model checking, fault injection and runtime
  verification \\
Correctness of goal expression & Is the specification the machine executes really the goal
  the human wants? &
  \textbf{Not derivable automatically from a program-correctness proof} \\
\bottomrule
\end{tabular}
\end{table}

The last row of Table \ref{tab:five-correctness} is the crux of the whole problem: \emph{a
program may implement a wrong specification correctly}. A compiler that is entirely correct
can compile a wrong program into machine code entirely correctly; a formal proof is equally
unable to rescue a specification that omits a critical requirement. The historical
statement of this distinction is the separation of verification from validation in software
requirements engineering \cite{boehm1984}.

\section{The Verifiable Refinement Chain}\label{sec:chain}

\begin{definition}[Verifiable refinement chain]\label{def:chain}
A \term{verifiable refinement chain} consists of seven stages:
\begin{enumerate}[label=(G\arabic*)]
\item \term{Purpose and intent}: what is to be achieved; which outcomes count as success;
      which side effects are unacceptable; how to trade off conflicting goals.
\item \term{System and environment model}: a representation of runtime states, events,
      operations, resources, operating-system interfaces, the external environment and
      observable behaviour.
\item \term{Formal behaviour contract}: permitted behaviour, forbidden behaviour,
      invariants that must hold, progress conditions, resource budgets and environmental
      assumptions, as in Definition \ref{def:contract}.
\item \term{Realisability decision}: check whether the contract is self-contradictory, and
      whether the goal is genuinely achievable under the specified resources, hardware,
      operating system and environmental conditions.
\item \term{Implementation synthesis and search}: generators, search algorithms or other
      synthesisers produce candidate implementations. A candidate implementation is not
      trusted in advance.
\item \term{Correctness proof from abstract behaviour to binary}: check the candidate
      implementation together with its compilation, linking and other transformations, and
      prove that the generated machine code satisfies the behaviour contract in the
      specified environment.
\item \term{Attested deployment and actual execution}: ascertain that what is running is
      the verified binary, execute it in the specified environment, and monitor, constrain
      and handle failures for the conditions that cannot be guaranteed statically.
\end{enumerate}
\end{definition}

The significance of this chain lies not in adding seven process steps but in \emph{assigning
a different mathematical task to each}. The model (G2) should not be a natural-language
description but an interpretable model of system state, behaviour and environment; the
realisability decision (G4) asks whether the requirements are satisfiable; synthesis (G5)
searches for possible implementations; (G6) establishes a proof from the implementation to
the required behaviour; and (G7) ensures that the object actually executed is the object the
proof is about.

\begin{principle}[What crosses a boundary]\label{pr:boundary}
At every boundary what crosses is not only an artefact but also \emph{what that artefact
stands for}, and \emph{the evidence for why it may be accepted}.
\end{principle}

\begin{remark}[The entry point that cannot be mechanically replaced]
Verification results and counterexamples can feed back into (G2), (G3) and (G4); only the
attestation of intent in (G1) cannot be replaced by a general code-correctness proof. The
reason is given in \S\ref{sec:intent-gap}.
\end{remark}

\section{Transition Edges and Evidence}\label{sec:five-edges}

The seven stages of Definition \ref{def:chain} compress into five transition edges, each
with its typical evidence and its critical gap. Table \ref{tab:edges} states them.

\begin{table}[htbp]
\centering
\footnotesize
\caption{Five transition edges, typical evidence and critical gaps}\label{tab:edges}
\begin{tabular}{L{1.5cm}L{3.4cm}L{4.6cm}L{4.9cm}}
\toprule
\textbf{Edge} & \textbf{Input $\to$ output} & \textbf{Typical proof/evidence} &
  \textbf{Key assumptions and principal gaps} \\
\midrule
I$\to$R & Natural language/dialogue $\to$ structured requirements, domain facts and
  items requiring clarification &
  Stakeholder attestation, links to requirement sources, examples and counterexamples,
  conflict reports &
  Intent has no automatically given mathematical truth value; external facts and conflicts
  of interest cannot be settled by type checking alone \\
R$\to$S & Requirements $\to$ formal properties/contracts &
  Traceability matrices, requirement-level tests, consistency/satisfiability/realisability
  analysis, vacuity and mutation analysis &
  A formal mapping must preserve \emph{prohibitions and positive examples} alike; an
  omitted requirement cannot be derived from internal consistency alone \\
S$\to$P & Specification $\to$ program/system candidates &
  Synthesis algorithms, SMT, model checking, Hoare/type/dependent-type proofs,
  counterexample-guided inductive synthesis &
  Finite DSLs, logic fragments and search budgets; over-strong premises and vacuous
  specifications must be prevented \\
P$\to$C$\to$B & Source program $\to$ stage-by-stage intermediate representations $\to$
  machine code/deployable artefact &
  Verified-compiler theorems, per-translation verification/attestation, proof-carrying code &
  Source and target semantics must be defined; linking, loading, assembly, the runtime,
  third-party libraries and hardware are not necessarily within the proof's scope \\
B$\to$E & Binary and configuration $\to$ behaviour in the real environment &
  Deployment integrity, system verification, runtime monitoring, isolation and capability
  mechanisms, performance and scheduling analysis &
  The model must be aligned with the real CPU, operating system, network, scheduling,
  failure and service assumptions \\
\bottomrule
\end{tabular}
\end{table}

\begin{remark}[Heterogeneity of the relations $R_i$]
By Remark \ref{rmk:source}, the interpretation relation on different transition edges need
not be of the same kind. For I$\to$R and R$\to$S the available ``relation'' may be only a
traceability mapping rather than a formal relation \cite{mucha2024,mokos2020}; for
P$\to$C$\to$B it may be a compiler's semantics-preservation theorem
\cite{leroy2009,cakeml2019}, a single translation verification \cite{necula2000} or a
translation attestation \cite{krijnen2024}; for B$\to$E it may be a decision by runtime
monitoring \cite{taleb2023}. Forcing them into one kind of relation is a mistake.
\end{remark}

\begin{remark}[Existing engineering on each edge]
Table \ref{tab:edges} is not a purely speculative scheme: verified compilation
\cite{leroy2009} and verified compiler back ends \cite{cakeml2019} on P$\to$C$\to$B,
kernel-level refinement proofs \cite{klein2009}, and proof-carrying code \cite{necula1997}
have already achieved substantial progress for specific languages, compilers, kernels and
systems. What remains genuinely unresolved is how to extend these high-assurance local
results to a wider range of targets and a greater variety of runtime environments, and to
systems whose implementations may be generated arbitrarily or adjusted dynamically.
\end{remark}

\section{Composition Conditions and Failure Modes}\label{sec:composition}

Proposition \ref{prop:soundness} already gives sufficient conditions for composition. Three
concrete counterexamples of composition failure are added here; their role is to show that
``locally correct'' does not entail ``globally correct''.

\begin{example}[Three composition failures]\label{ex:compose-fail}
\begin{enumerate}[label=(\alph*)]
\item The source program is semantically correct, but the linker selects the wrong library.
\item Each service satisfies its interface contract on its own, but the contracts'
      assumptions about network retry are mutually incompatible.
\item The functional proof holds, but the latency budget does not cover cache contention.
\end{enumerate}
The common structure of the three is that one of conditions (a) or (b) of
Proposition \ref{prop:soundness} is violated.
\end{example}

Read through the contract algebra (Definition \ref{def:operad},
Proposition \ref{prop:operad-compose}), the three counterexamples are no longer mutually
unrelated; they fall into three classes of failure:

\begin{corollary}[Common structure of the three counterexamples]\label{cor:three-failures}
\begin{enumerate}[label=(\alph*)]
\item Linker selects the wrong library: \emph{colour mismatch}, the colour of the linked
      implementation being unequal to the colour declared at the link point;
\item Contracts' assumptions about network retry are mutually incompatible: \emph{a cycle
      appears in the interface graph}, retry turning an acyclic call graph into a cyclic
      one while the colouring is not re-verified after the cycle appears
      (Remark \ref{rmk:feedback});
\item The latency budget does not cover cache contention: \emph{a missing contract
      component}, the resource component $R$ never entering the colour, so that cost is
      treated as an implementation detail.
\end{enumerate}
The three are, respectively, a colour whose extension is too small, a changed interface
graph, and a missing contract component.
\end{corollary}

The conclusion is that composition rules must explicitly incorporate these interface
conditions together with \emph{non-functional conditions}.

\begin{remark}[The holistic character of assurance]\label{rmk:duhem}
Corollary \ref{cor:three-failures}(c) and Remark \ref{rmk:monitor-limit} together show that
no verification conclusion actually asserts ``property $\varphi$ holds''; what it asserts is
the conjunction
\[
\varphi \;\land\; A_1 \;\land\; A_2 \;\land \cdots \land A_k,
\]
where the $A_i$ are the full set of environmental assumptions. The number $k$ cannot be
enumerated in advance, and an assumption $A_i$ becomes visible only when it is actually
violated.
This is the form taken by the Duhem--Quine thesis \cite{quine1951} in software verification:
there is no decisive test of a single proposition, only a test of the whole conjunction. Its
immediate consequence is that the strength of ``the tests pass'' is strictly limited by the
\emph{set of assumptions that were not enumerated}. This also explains why
Table \ref{tab:evidence} must attach a ``what must not be claimed'' list to every class of
evidence: that list is the known $A_i$.
\end{remark}

\section{The Intent Gap}\label{sec:intent-gap}

This is the section that requires the most care.

\begin{problem}[The intent gap]
How is an intent $I$, first expressed by a human in natural language, guaranteed to be
transformed into a correct formal specification $\Phi$? The hope is that $I \equiv \Phi$,
where the equivalence means that ``the behaviour $\Phi$ requires fully reflects the human's
real goal''.
\end{problem}

A popular argument holds that, by G\"odel's incompleteness theorem and Rice's theorem, the
gap between intent and specification is \emph{principled} and \emph{can never be
mechanically closed}. The argument does not hold.

\begin{proposition}[Why this paper does not accept that argument]\label{prop:no-godel}
G\"odel's incompleteness theorem concerns the limits of \emph{formal provability} in formal
systems satisfying certain conditions; it does not directly state that a machine cannot
decide whether a requirements specification matches a particular person's real intent.
Rice's theorem \cite{rice1953} states that non-trivial semantic properties of general
programs are undecidable; that is a fundamental limit on general static analysis, but it is
not a theorem about the relation between a person's mental intent and a formal
specification. Deriving ``intent can never be fully formalised'' from it requires additional
philosophical and cognitive assumptions. Turing's original result \cite{turing1936}
likewise applies only to explicitly encoded computational problems and general algorithmic
claims.
\end{proposition}

Does the intent gap exist, then? It does, but the grounds are \emph{epistemic}, not
\emph{logical}.

\begin{proposition}[Three epistemic grounds for the intent gap]\label{prop:three-grounds}
\begin{enumerate}[label=(\roman*)]
\item \term{Requirements are generative, not enumerative}. A person cannot list their
      requirements completely before seeing the system's behaviour. The requirement ``do
      not leak privacy'' may not surface at all until the system is seen actually
      collecting data. This is not logical undecidability but a constructive limitation of
      human cognition: \emph{preferences are activated by the environment, not present
      beforehand as a complete list}.
\item \term{A specification is lossy compression of intent}. Intent lives in a causal,
      embodied, evolutionarily shaped mind; a specification is a finite symbolic object.
      Lossy compression means that information is necessarily lost, and which parts are
      lost cannot be known in advance. This is a phenomenon of information theory, not of
      metamathematics.
\item \term{The distinction between verification and validation has not been overturned}.
      Verification answers ``does the implementation conform to the specification'';
      validation answers ``does the system meet the real need''. The distinction exists and
      persists precisely because the latter cannot be mechanically reduced to the former.
\end{enumerate}
\end{proposition}

Combining Proposition \ref{prop:no-godel} with \ref{prop:three-grounds} gives the paper's
precise statement of the problem:

\begin{principle}[The precise statement of the intent gap]\label{pr:gap}
The specification-transformation process is not a step that can be mechanically guaranteed
\emph{in general}; the reason is not logical undecidability but that human preferences are
dynamic, situational, and not fully knowable before they are encountered.
\end{principle}

This correction does not weaken the claim; it replaces the load-bearing wall. It also shows
that machines can help clarify requirements, propose counterexamples, generate scenarios,
find contradictions and check a specification for internal consistency, whereas which
specification is finally adopted, and whether it fully expresses the goal, still requires an
identifiable accountable party.

\begin{remark}[Three questions, not one, and what automation can do for each]\label{rmk:three-questions}
The intent gap is often discussed as a single obstacle. It is three, they need different
instruments, and separating them makes clear that the framework denies mechanisability to a smaller
region than it is usually taken to.

\begin{enumerate}[label=(\alph*)]
\item \term{Semantic expression}: given requirements that have in fact been settled, can they be
      written as a formal specification? In bounded domains this is largely automatable, and the
      instruments are the type system, the theorem prover and the specification-synthesis methods
      discussed in Chapter \ref{ch:literature}. Nothing in Proposition \ref{prop:three-grounds}
      denies this, and treating it as denied would be an over-reading.
\item \term{Requirement discovery}: has the interested party found and expressed everything that
      matters? This is Proposition \ref{prop:three-grounds}(i), and its remedy is interaction:
      counterexamples, scenarios, and the probing loop of Chapter \ref{ch:comprehension}. Partially
      automatable, in the sense that machines can generate the probes; not automatable in the sense
      that the answers must come from someone with something at stake.
\item \term{Normative choice}: when requirements conflict---latency against safety, privacy against
      analytical power, cost against performance---who has the standing to decide? This is not an
      epistemic question at all; it is the authority axis of \S\ref{sec:authority}, and the
      instruments are procedural: a designated accountable body and a record of the decision.
\end{enumerate}
Reading the gap this way also corrects a common over-statement in the opposite direction. Because
(a) is largely mechanisable, ``some intent can be precisely defined and automatically checked'' is
true; what cannot be mechanised is the residue---preferences not yet activated, conflicts not yet
surfaced, and decisions nobody has the standing to make alone. The precise claim this paper defends
is therefore narrower than ``intent cannot be formalised'': \emph{a formal proof of implementation
correctness cannot by itself establish whether the specification serves the goals of all relevant
parties}, and the open-world remainder must be handled by an independent validation process. That
process may use automated tools; it may not treat a tool's internal consistency as evidence that the
goal was right.
\end{remark}

\begin{remark}[The intent gap is the frame problem]\label{rmk:frame}
The three grounds of Proposition \ref{prop:three-grounds} are not three independent
findings; they point jointly to an older problem. To formalise an intent one must first
decide which facts are \emph{relevant}; and deciding relevance requires already knowing
which conclusions are to be drawn, which is exactly the frame problem: without knowing the
answer in advance, one cannot enumerate the relevant premises.

The intent gap is therefore not a defect of formal methods but an instance of the frame
problem. The identification has two consequences:
\begin{enumerate}[label=(\alph*)]
\item It explains why stronger formal tools do not solve the problem: tools improve
      \emph{derivation}, whereas the problem lies in \emph{the enumeration of premises};
\item It explains why \emph{interaction} is the only effective means. The frame problem is
      unsolved in general and is evaded only in systems that possess a perception--action
      loop: an agent need not enumerate relevance in advance, because the environment gives
      feedback on its actions. The probing loop of Chapter \ref{ch:comprehension} and the
      traceable refinement of this section both move that evasion mechanism up to the
      requirements layer.
\end{enumerate}
Interactive refinement is therefore forced, not selected from among candidate approaches.
\end{remark}

\begin{remark}[The oblivious set is where the frame problem is paid for, level by level]\label{rmk:frame-per-level}
Remark \ref{rmk:frame} identifies the intent gap with the frame problem, and stops there. The
identification can be pushed one step further, and doing so joins two parts of the paper that have so
far been separate.

The frame problem, in the form used here, is the difficulty of \emph{deciding relevance}: of saying
which facts bear on a conclusion before knowing what is to be concluded. An oblivious set
(Definition \ref{def:oblivious}) is a relevance judgement of exactly that shape, made explicit and
local: $\Obl(\Lambda_i)$ names the differences that $\Lambda_i$ declares not to bear on it. The
relation between the two is therefore not analogy but scale---and it is worth being precise about
what is and is not obtained.

What is obtained is a \emph{localisation}. A frame problem posed for a system as a whole is not
smaller for being posed for one of its levels; but it is bounded by that level's declared
observation map, and it is discharged in a currency the level can state. The global question, which
facts are relevant to the system's purpose, still has no general solution, and nothing in this paper
claims one. What the level structure buys is that the global question is never asked all at once:
each level answers its own bounded version and publishes the answer.

What is obtained second is \emph{testability}. A relevance judgement that is merely held cannot be
criticised; one that is written down can. Section \ref{sec:neardecomp} gives the two ways a declared
judgement can be wrong---the interface may be insufficient for external behaviour, or the observation
map may be incomplete with respect to what is observable in fact. Both are checks, not opinions. In
this sense the oblivious set is the frame problem's engineering form: not a solution to it, but a
place to pay for it where the payment can be audited.

Two consequences follow, and they are the reason the remark is here rather than in Chapter
\ref{ch:comprehension}. First, it explains why the descent protocol of Principle \ref{pr:descent} has
the stopping rule it has: one descends until the invariant governing the change becomes visible,
because that is the depth at which the relevant oblivious sets have been consulted, and descending
further buys relevance that no obligation depends on. Second, it makes the intent layer's
difficulty legible in the same terms as a module boundary's: the intent gap is a frame problem at
$\Lk{\infty}$, which is exactly the level at which no oblivious set can be declared for the machine,
because declaring one would already be deciding the purpose.
\end{remark}

\section{The Asymmetry Between Verification and Generation}\label{sec:asymmetry}

A second formulation requiring correction is the description of ``grading is easier than
sitting the exam'' as ``the universal dividend of $\mathrm{P}$ versus $\mathrm{NP}$''. The
description does not hold: whether $\mathrm{P}$ equals $\mathrm{NP}$ remains open, and it is
not known in general whether finding a solution is necessarily harder than checking one.

This correction must be distinguished from a second point: the open status of
$\mathrm{P}$ versus $\mathrm{NP}$ does \emph{not} make the phenomenon ``verification is
cheaper than generation'' fictitious.

\begin{proposition}[Structural asymmetry]\label{prop:asym-struct}
The structural asymmetry of polynomial-time certificate checking is \emph{the definition of}
the class $\mathrm{NP}$ itself: that a solution can be verified efficiently is a basic fact
about computation, and it does not depend on whether $\mathrm{P} \neq \mathrm{NP}$ has been
proved.
\end{proposition}

\begin{proof}[Argument]
The definition of $\mathrm{NP}$ is ``there exists a polynomial-time verifier $V$ and a
certificate $w$ of polynomial length such that
$x \in L \iff \exists w: V(x,w) = 1$''. The definition itself asserts the existence and the
efficiency of verification, and asserts nothing about the complexity of finding $w$. The
effectiveness of the verifier architecture therefore rests on the definition itself and is
independent of the eventual answer to $\mathrm{P}$ versus $\mathrm{NP}$.
\end{proof}

This cost scale indicates the precise boundary of ``verification is cheap''. The limit is
not the open status of $\mathrm{P}$ versus $\mathrm{NP}$ but \emph{certificate length}:
checking a proof can be done in polynomial time, while the proof itself can be
astronomically long. Hence:

\begin{principle}[The two components of verification cost]\label{pr:two-costs}
``Verification is cheaper than generation'' must be decomposed into two independent
propositions:
\begin{enumerate}[label=(\alph*)]
\item \term{Checking time is cheap}: given a certificate, the decision is efficient
      (Proposition \ref{prop:asym-struct});
\item \term{Certificate size is cheap}: the certificate itself is short enough to be
      stored, transmitted and re-checked.
\end{enumerate}
(a) is guaranteed by the definition of $\mathrm{NP}$; (b) is \emph{not}, and in many
practical systems it is the real bottleneck. A verification regime that satisfies only (a)
is unusable in engineering, because no one can re-check a certificate they cannot finish
reading.
\end{principle}

Principle \ref{pr:two-costs}(b) is measurable: for a given prover and class of properties it
suffices to measure the distribution of certificate sizes. It also explains why a ``small
kernel plus large proof'' architecture must in practice always be paired with a
\emph{readable} proof-summary layer (invariants, prose statements of contracts,
counterexamples) rather than supplying only proof terms.

\begin{example}[Historical cost scale of formal verification]\label{ex:sel4-cost}
The initial functional-correctness proof of the microkernel seL4 took about 11 person-years
(excluding investment in reusable tools and frameworks); the proof scripts of the time ran
to about 200{,}000 lines and verified about 8{,}500 lines of C code \cite{klein2009}. That
scale indicates the historical cost structure of formal verification. AI may substantially
reduce the cost of writing proofs, finding invariants and repairing proofs, but
specification design, model selection, proof architecture and the trusted base remain real
work.
\end{example}

\section{Grades of Evidence}\label{sec:evidence-levels}

Different kinds of evidence provide support of different strength, and conflating them is
the most common error in this area. Table \ref{tab:evidence} gives the grades of evidence
together with the corresponding ``what must not be claimed'' list.

\begin{table}[htbp]
\centering
\footnotesize
\caption{Grades of evidence}\label{tab:evidence}
\begin{tabular}{L{2.6cm}L{5.6cm}L{6.2cm}}
\toprule
\textbf{Type of evidence} & \textbf{Conclusions it can support} &
  \textbf{Conclusions that must not be claimed} \\
\midrule
Machine-checked formal proof & For the stated proposition, holds under the given logic,
  definitions, axioms, libraries and checker assumptions &
  Does not prove that the proposition fully expresses the intent; does not cover unmodelled
  environments; does not by itself establish latency or energy consumption \\
Verified-compiler theorem/translation certificate & Between specific source and target
  semantics, the defined preservation or refinement relation holds, or a single translation
  result is verified &
  Does not mean that the linker, operating system, hardware and external services are
  verified \\
Static analysis and abstract interpretation & Proves or conservatively estimates the
  properties of the analysed program, relative to a sound approximation &
  Does not guarantee precision or the absence of false alarms; is not a measurement of
  actual performance \\
Testing and mutation testing & Finds counterexamples, assesses the specified behaviour on
  the listed tests, and provides evidence over a finite sample &
  Does not prove that all inputs and all executions are safe; a pass rate is not a general
  proof of semantic equivalence of specifications \\
Statistical performance evidence & An empirical distribution under specified data, sampling
  procedure, hardware and confidence method &
  Is not a hard-real-time worst-case execution time; cannot be extrapolated unconditionally
  to a deployment whose distribution has changed \\
Runtime monitoring & Gives an online decision about the observed traces under the monitor's
  semantics &
  If observation is incomplete, ``not detected'' must not be stated as ``all executions
  satisfy'' \\
\bottomrule
\end{tabular}
\end{table}

\begin{principle}[Three states that must be distinguished]\label{pr:three-states}
``A violation was detected'', ``no violation has yet been observed'' and ``all executions
are proved safe'' are three different states, and none may be substituted for another.
\end{principle}

\begin{definition}[Artefact evidence list]\label{def:evidence-package}
Every artefact should carry ten items of evidence, listed in
Table \ref{tab:evidence-package}. The purpose of the list is to make ``may this binary be
trusted'' an auditable question rather than a question about the identity of its author.
\end{definition}

\begin{table}[htbp]
\centering
\small
\caption{Artefact evidence list}\label{tab:evidence-package}
\begin{tabular}{L{0.8cm}L{13.4cm}}
\toprule
\textbf{No.} & \textbf{Evidence item} \\
\midrule
1 & Specification version and requirement traceability identifiers \\
2 & Formal specification and environmental assumptions \\
3 & Hash of the candidate source code \\
4 & Compiler, flags, target architecture and dependency versions \\
5 & Proof and checker versions \\
6 & Resource-analysis model and results \\
7 & Test and monitoring strategy \\
8 & Trusted computing base inventory \\
9 & \textbf{Known uncovered behaviour} \\
10 & \textbf{Re-verification rules} (which changes trigger recomputation of the affected
  evidence) \\
\bottomrule
\end{tabular}
\end{table}

Items 9 and 10 are the two most easily omitted and the most important: the former makes the
boundary of the guarantee explicit, the latter turns assurance from a one-off check into a
lifecycle mechanism.

\section{Side Effects and Constraint Mechanisms}\label{sec:effects}

The example in Problem \ref{prob:trust} mixes side effects of several kinds. Their
mechanical mechanisms become visible only when the kinds are separated.

\begin{definition}[Four classes of side effect]\label{def:effects}
See Table \ref{tab:effects}.
\end{definition}

\begin{table}[htbp]
\centering
\small
\caption{Classification of side effects and the corresponding mechanical mechanisms}\label{tab:effects}
\begin{tabular}{L{1.8cm}L{4.6cm}L{5.0cm}L{3.0cm}}
\toprule
\textbf{Class} & \textbf{Example} & \textbf{Mechanical mechanism} &
  \textbf{Point of verification} \\
\midrule
Functional & A generated function modifies global state it must not modify &
  Type systems, effect types (which make ``what will be touched'' part of the signature),
  ownership systems &
  Compile time, at zero runtime cost \\
Resource & 1\,ns becomes 1\,$\mu$s; memory blows up &
  Fuel/gas metering, worst-case execution time analysis, quotas &
  Both compile time $+$ runtime \\
Security & Unauthorised access, sandbox escape &
  Capability models, sandboxes, bytecode verifiers &
  At admission $+$ at runtime \\
Semantic drift & The specification itself is wrong &
  \textbf{No mechanical mechanism} & Human attestation and consensus \\
\bottomrule
\end{tabular}
\end{table}

The first three classes are already covered by mechanical mechanisms, and are arranged in
depth (proof $>$ types $>$ testing $>$ metering $>$ isolation $>$ monitoring, each layer
having a simpler trust base than the layer above it). Only the fourth class, a specification
that departs from the intent, always requires humans.

\begin{remark}[Assurance scope of three engineering prototypes]\label{rmk:three-prototypes}
Three common engineering mechanisms cover different kinds of property and must not be
advertised side by side:
\begin{enumerate}[label=(\alph*)]
\item A browser sandbox guarantees \emph{execution isolation and a permission boundary}, not
      functional correctness. It is the largest untrusted-code runtime on earth: every day
      it executes hundreds of millions of pieces of code of unknown provenance that have
      never been reviewed; it never reviews, it isolates.
\item A bytecode verifier verifies \emph{the safety properties its rules cover} (no
      out-of-bounds access, bounded loops, no dangling pointers), not the correctness of
      arbitrary business logic. A deterministic verifier of a few tens of thousands of
      lines adjudicates the admission of millions of lines of free code; this is precisely
      an instance of ``compiler-like assurance that does not depend on anyone's will''.
\item A resource-metering mechanism guarantees \emph{constraints under an abstract resource
      accounting}, not real wall-clock time. It turns a performance side effect from
      something ``a human must notice'' into ``a budget violation automatically refuses
      execution'', but it need not catch a degradation from 1\,ns to 1\,$\mu$s, because it
      meters abstract operation costs rather than actual hardware latency. Real-time
      performance assurance requires timing analysis and a restricted runtime environment,
      a completely different engineering discipline.
\end{enumerate}
\end{remark}

\section{Runtime Assurance}\label{sec:runtime}

Formal proofs depend on environmental assumptions, and the actual environment may change.
Beyond static proof, therefore, the executing system must remain safe when an assumption
fails.

\begin{principle}[Layered defence]\label{pr:layered-defense}
An ideal architecture combines static proof (for known deterministic behaviour) with dynamic
monitoring (for unknown or non-deterministic behaviour), forming a layered defence: whatever
can be proved is proved first; residual risk is detected by independent monitoring; serious
violations have their impact limited by isolation or rollback.
\end{principle}

\begin{remark}[The monitor's own boundary]\label{rmk:monitor-limit}
Even a runtime monitor that is itself correct does not thereby possess the complete facts of
an execution. Missing observation traces, ambiguous events, unreliable transport and
concurrency that cannot be seen all limit its conclusions. The observation mechanism of the
monitoring system must therefore itself be included in the assurance chain.
\end{remark}

\begin{remark}[Four conditions for contract composition]\label{rmk:contract-conds}
The composition of local properties depends on four things: the environment genuinely
satisfies the subcomponent assumptions; the connection topology is well posed; the feedback
condition is checkable; and the observation and time semantics are compatible. ``Individually
correct'' for components and ``globally correct'' for the system are not interchangeable.
\end{remark}

\begin{principle}[The conditional form of an invariant]\label{pr:ag}
A usable assume--guarantee contract should state simultaneously: the environmental
conditions it depends on, the behaviour it guarantees outwards, the side effects it may
produce, the resources it must consume and occupy, and the constraints that must hold when
it is composed with other subsystems.
\end{principle}

\section{Artefact Binding}\label{sec:fourth}

The last conjunct of the assurance condition \eqref{eq:assurance}, namely that the object
actually executed is the object verified, is frequently omitted from existing discussions.
It is not a marginal detail: even if a verifier correctly accepts a binary, the preceding
verification conclusion does not transfer automatically if the deployment system uses a
different file, or if the loading environment changes a critical premise. Supply-chain
attacks aim precisely at this crack.

The complete trust chain must therefore include content-addressed build outputs, a
cryptographic binding between the verification evidence and the binary, and re-verification
at deployment. In the notation of \eqref{eq:binding},
$\mathrm{digest}(B_{\text{deployed}}) = \mathrm{digest}(B_{\text{verified}})$
and $E_{\text{deployed}} \models A_{\text{claimed}}$.

\begin{principle}[A three-way distinction for reproducible builds]\label{pr:reproducible}
Three properties of different strength must be distinguished:
\begin{enumerate}[label=(\alph*)]
\item \term{Reproducible}: the same inputs produce a byte-for-byte identical output;
\item \term{Verifiable}: an independent third party can rebuild from the inputs and compare;
\item \term{Verified}: a proof that a particular binary satisfies its declared properties.
\end{enumerate}
(a) and (b) are preconditions for (c) to take effect; (c) does not imply (a). The three are
not interchangeable.
\end{principle}

\section{The Reliability Bound of the Consensus Layer}\label{sec:consensus}

\S\ref{sec:trust} showed that the intent layer admits no mechanical solution. What does it
admit? This section gives a precise answer and proves a counter-intuitive result.

\begin{theorem}[Consensus reliability bound]\label{th:condorcet}
Let $n$ verifiers judge independently the proposition ``whether a specification faithfully
expresses the intent'', each erring with probability at most $p < 1/2$. Then the probability
that majority voting errs satisfies
\begin{equation}\label{eq:chernoff}
\Pr\bigl[\text{majority errs}\bigr] \;\le\; \exp\Bigl(-2n\Bigl(\tfrac{1}{2} - p\Bigr)^{2}\Bigr).
\end{equation}
\end{theorem}

\begin{proof}
Let $X$ be the number of verifiers who err, $X = \sum_{i=1}^{n} X_i$, with the $X_i$
independent and $\Pr[X_i = 1] \le p$, so $\mathbb{E}[X] \le np$. A majority errs exactly
when $X \ge n/2$. By Hoeffding's inequality \cite{hoeffding1963}, for $t > 0$,
$\Pr[X - \mathbb{E}[X] \ge t] \le \exp(-2t^{2}/n)$. Taking
$t = n(1/2 - p)$ (so that $X \ge n/2$ implies $X - \mathbb{E}[X] \ge t$) gives
\[
\Pr\Bigl[X \ge \tfrac{n}{2}\Bigr] \;\le\;
\exp\Bigl(-\frac{2n^{2}(1/2 - p)^{2}}{n}\Bigr) \;=\;
\exp\Bigl(-2n\bigl(\tfrac{1}{2} - p\bigr)^{2}\Bigr). \qedhere
\]
\end{proof}

This result is the quantitative form of the jury theorem \cite{condorcet1785}. It relies on
independence, however, and independence is a fragile assumption; its failure has an equally
precise form.

\begin{proposition}[Bound dominated by a correlated error mode]\label{th:correlated}
Suppose that the errors of $n$ verifiers have the following structure: with probability $q$
a \emph{common error mode} $C$ occurs (for instance a shared training-data bias, a shared
domain misunderstanding or a shared incentive), in which case all verifiers err
simultaneously; with probability $1-q$ the verifiers err independently, with error rate
$p' < 1/2$. Then
\begin{equation}\label{eq:correlated}
\Pr\bigl[\text{majority errs}\bigr] \;\ge\; q .
\end{equation}
In particular, when $q$ is bounded below by a positive constant, increasing $n$ cannot push
the probability that the majority errs below $q$.
\end{proposition}

\begin{proof}
When $C = 1$ all $n$ verifiers err, so the majority necessarily errs. Hence
$\Pr[\text{majority errs}] \ge \Pr[C = 1] = q$. Taking the limit in $n$ does not change
this lower bound.
\end{proof}

\begin{corollary}[Diversity, not number]\label{cor:diversity}
By Proposition \ref{th:correlated}, the governing variable of consensus reliability is
\emph{the independence of errors} (that is, the size of $q$), not the number of verifiers.
The central goal of consensus design is therefore to reduce $q$: admit participants with
different backgrounds, models with different training provenance, reviewers with different
interests, and explicitly designed \emph{adversarial} roles.
\end{corollary}

\begin{remark}[Use of this result]\label{rmk:diversity-use}
Corollary \ref{cor:diversity} yields a criterion directly usable in engineering judgement:
when someone proposes ``adding a few more agents to review this'', the right question is
``are the new reviewers independent of the existing ones in their error modes''. If they are
not, increasing the number increases cost without increasing reliability. This also explains
why ``reviewing with another model from the same source'' does not constitute a guarantee
structurally.
\end{remark}

\begin{remark}[What the consensus model is, and what it is not]\label{rmk:consensus-status}
Theorem \ref{th:condorcet}, Proposition \ref{th:correlated} and the design-effect approximation
below are statements \emph{within} stated probability models. They should not be read as predictions
of the reliability of a real review process, and the gap between model and practice is worth being
explicit about, because it is where the framework is most likely to be misused.

Three caveats. First, the models assume a common proposition to be judged, independent error sources
modulo $q$, and a fixed individual competence. Real attestation exercises violate all three: what
counts as a faithful specification is itself under negotiation, attestors influence one another by
talking, and competence varies with the clause being judged. The literature on jury theorems is
largely a study of how fragile these premises are---in particular, independence and competence are
not simultaneously justifiable under any single notion of probability \cite{dietrich2008}---and a
shared body of evidence is the most natural mechanism by which apparent independence fails, since it
makes the attestors' errors share a common cause \cite{dietrich2004}. The second of those results is
the formal counterpart of what Principle \ref{pr:q-decomposition} calls information isomorphism.

Second, a single parameter $q$ does not characterise dependence in general. Two attestors can be
correlated in one clause and independent in another, correlated through shared training in one
subdomain and through shared incentives in another, and a one-dimensional summary averages over all
of it. $q$ is the right parameter for the bound and a lossy summary of the phenomenon.

Third, and most consequentially, $p$ and $q$ are not estimable from data unless ``faithful to the
intent'' has a stable standard of evaluation, a designated adjudicator and a specification of what
counts as an error. Where those exist---bounded domains with observable acceptance criteria---the
quantities can be estimated from outcomes, which is the case Hypotheses $H_8$ and $H_{11}$ address.
Where they do not, the bound remains a conditional statement and no quantity of attestation data
supplies the missing standard. The honest summary is that the model tells you what to measure if the
measurement is possible, and tells you the shape of the answer if it is.
\end{remark}

\begin{remark}[Partial correlation: the design effect]\label{rmk:design-effect}
Theorem \ref{th:condorcet} and Proposition \ref{th:correlated} describe two extremes: independent
errors, giving a bound that improves exponentially in $n$, and a common error mode, giving a floor
$q$ that no $n$ can breach. A real attestation network sits between them, with errors partially
correlated in pairs. The standard statistical instrument applies, and is worth writing down because
it makes Corollary \ref{cor:diversity} quantitative rather than rhetorical.

Let $\bar\rho$ be the mean pairwise correlation of attestor errors. The customary
\emph{design-effect} approximation \cite{kish1965} replaces the nominal attestor count by an
effective one,
\begin{equation}\label{eq:design-effect}
n_{\mathrm{eff}} \;=\; \frac{n}{1 + (n-1)\bar\rho},
\end{equation}
and one reads the bound of Theorem \ref{th:condorcet} at $n_{\mathrm{eff}}$ in place of $n$. Two
consequences follow.
\begin{enumerate}[label=(\alph*)]
\item As soon as $n\bar\rho \gg 1$, \eqref{eq:design-effect} gives
      $n_{\mathrm{eff}} \to 1/\bar\rho$, \emph{independent of $n$}. Beyond that point additional
      attestors buy nothing; the ceiling is set by the correlation alone.
\item Lowering $\bar\rho$ raises that ceiling, and the trade is explicit: once $n\bar\rho$ is large,
      halving $\bar\rho$ roughly doubles the achievable $n_{\mathrm{eff}}$. Diversity is a budget
      line, not a slogan.
\end{enumerate}
The relation to Proposition \ref{th:correlated} is that $q$ is the degenerate case
$\bar\rho \to 1$: with complete correlation $n_{\mathrm{eff}} \to 1$ and \eqref{eq:design-effect}
reproduces the floor.

Two caveats belong with the remark. First, \eqref{eq:design-effect} is the standard approximation
for the variance of a mean under exchangeable correlation, not a theorem about majority error; a
rigorous bound for the intermediate regime needs an explicit correlation model, and none is claimed
here. Second---and this is what makes the approximation usable---$\bar\rho$ is \emph{estimable} from
data the framework already collects: the historical pairwise correlation of attestor errors, settled
against subsequently observed defects and rework. This does not answer the open problem
\ref{op:q} of measuring $q$; it supplies a companion quantity that is measurable, while $q$ remains
the limit that is not.
\end{remark}

\subsection{Why the two layers cannot be merged, and what that forces}\label{sec:enrichment}

Definition \ref{def:two-layer} asserts that the chain has a mechanical layer and a consensus layer
and that the two are stitched rather than fused. That assertion has a mathematical reason, and
stating it settles a design question that is otherwise answered by preference: why a single, more
capable verifier cannot replace the pair.

The reason is that the two layers do not merely differ in strength; they assign their guarantees to
\emph{different kinds of object}. A mechanical guarantee is a proof: an arrow either carries one or
it does not, and one proof may refine another. A consensus guarantee is a number: an unreliability
$\varepsilon$, with no question of one $\varepsilon$ refining another. Formally the two layers are
two \emph{enrichments} of the same underlying structure of stages and artefacts---hom-objects drawn
from the two-element order in the first case, from the Lawvere quantale
$V = ([0,\infty],\geq,+,0)$ in the second, where the composite of two stages has unreliability at
most the sum of theirs. That composition law is the union bound, and it deserves emphasis that it
needs no independence assumption: if stages $1$ and $2$ are each faithful with probability at least
$1-\varepsilon_i$, then both are faithful with probability at least $1-\varepsilon_1-\varepsilon_2$
whatever the dependence between them. This is worth contrasting with the bound of
Theorem \ref{th:condorcet}, whose exponential improvement requires independence and which the whole
of \S\ref{sec:consensus} exists to qualify.

The asymmetry between the two enrichments is not symmetric. There is a change of base from proof to
reliability---a discharged obligation has unreliability $0$, an undischarged one has
unreliability $\infty$---so a proof may legitimately be weakened into a reliability claim. The
reverse fails, and it fails for a reason that is two lines long. If
$G : ([0,\infty],\geq,+,0) \to (\{0,1\},\leq,\wedge,1)$ respects the monoidal structure, then for
every $\varepsilon$,
\begin{equation}\label{eq:doubling}
G(2\varepsilon) \;=\; G(\varepsilon+\varepsilon) \;=\; G(\varepsilon)\wedge G(\varepsilon)
\;=\; G(\varepsilon),
\end{equation}
so $G$ cannot distinguish $\varepsilon$ from $2^n\varepsilon$ for any $n$. Nothing is lost in this
computation by weakening it: functoriality alone gives $G(2\varepsilon) \leq G(\varepsilon)$ and the
comparison map gives the reverse inequality, so \eqref{eq:doubling} holds for lax monoidal $G$ as
well as strong ones. The form checked in the companion development is the strict one,
$G(a+b) = G(a)\wedge G(b) \Rightarrow G(2\varepsilon) = G(\varepsilon)$; the lax strengthening is the
two-line argument just given and is not itself checked.

The monoidal maps from $V$ to the two-element order are exactly three: the constant map to the top
element; the map $\varepsilon \mapsto [\varepsilon < \infty]$, which is the left adjoint of the
inclusion of the two-element order and is Lawvere's $\pi_0$; and the map
$\varepsilon \mapsto [\varepsilon = 0]$, which is the right adjoint \cite{lawvere1973}. \emph{Each of the three is constant on the finite positive values}---the first two take the top
element there and the third takes the bottom---so none of them separates two such values. That is
the consequence this chapter uses: \emph{no monoidal map separates two finite positive
unreliabilities}. A single threshold on $\varepsilon$ is therefore not merely a lossy summary; it is
not a change of base at all. With $\theta$ for the threshold and $G(\varepsilon) = [\varepsilon \leq
\theta]$, two stages at $\varepsilon = 0.6\,\theta$ each pass while their composite at
$1.2\,\theta$ does not, so $G$ fails \eqref{eq:doubling} and with it the composition law.

\begin{principle}[Reliability is budgeted, not gated]\label{pr:budget}
A chain of $k$ stages that must \emph{all} hold has composite unreliability at most
$\sum_{i=1}^{k}\varepsilon_i$. Gating each stage against a fixed threshold $\theta$ therefore gives a
composite bound of $k\theta$, which grows without limit in $k$; the bound must instead be
\emph{divided}, with each stage held to $\varepsilon_i \leq \Theta/k$ for a total budget $\Theta$.
\end{principle}

The practice the principle rules out is common enough to be worth stating numerically. Take
``$99\%$ reliable'' as the per-clause gate, $\varepsilon \leq 0.01$, and a specification with $k$
clauses each attested on its own---attested separately, not with independent errors, since the union
bound assumed no independence. The composite bound $1-k\varepsilon$ is $0.90$ at $k = 10$, $0.50$ at
$k = 50$, and \emph{zero at $k = 100$}---no information at all. Real specifications have hundreds of
clauses, so clause-wise gating at a fixed per-clause reliability does not merely weaken the
guarantee; past a hundred clauses it says nothing, while continuing to look like a guarantee. The
bound must be divided: to have any content at $k = 100$ the per-clause requirement is
$\varepsilon < 0.01$, strictly below $1/k$. This is the same shape as the budget line of Remark
\ref{rmk:design-effect}---which already says of diversity that it is a budget---applied to a
different quantity, and it is the shape the design-effect correction takes across stages rather than
within one: the design effect governs correlation \emph{among attestors of one stage}, the union
bound governs composition \emph{across stages}.

Two boundary conditions on the claim. First, it concerns the composition of \emph{reliability}; it
does not say that thresholds are useless, only that a threshold applied per stage is not a
compositional summary, and the correct composite statement is the $\varepsilon$-indexed family
$[\varepsilon \leq \theta]$ taken for each $\theta$ rather than one chosen $\theta$. That family is
how a quantitative statement is expressed in Boolean logic without losing its content, and the
device is standard in quantitative equational logic \cite{bacci2023}. The failure mode of a single
threshold has a standard name in the bisimulation-metric literature as well: approximate
bisimulations are not transitive in general, which is the abstract form of the
$0.6\theta + 0.6\theta > \theta$ computation \cite{gebler2013}.

Second, the principle bounds what a chain can claim, not what a single stage can achieve. The
$\varepsilon_i$ are properties of one attestation round given a design and a correlation structure;
the design effect is how they are estimated. The contribution here is therefore narrow and should be
stated narrowly: not a better bound, but the observation that the sum bound, which is often treated
as a crude fallback, is the composition law \emph{forced} by the enrichment, and that the
alternatives to it are not better bounds but a different kind of object. This is worth saying because
the opposite is the received view at exactly the point where the framework needs it. Rushby, in the
analysis of assurance cases closest to the present chapter, states the mismatch between probabilistic
evidence and deductive argument, gives the sum rule and the product rule, judges both ``too crude and
pessimistic'', and leaves the repair as an open problem \cite{rushby2015}. The refinement this
section offers is that the sum rule is not a crude approximation to be improved on---it is what
composition of unreliabilities \emph{is}---and that the apparent crudeness lies in the per-stage
$\varepsilon$, where the design-effect correction and the incentive analysis of \S\ref{sec:mechanism}
do their work, rather than in the rule that combines them.

One point of contact with the software-engineering literature is worth recording, because it is the
empirical counterpart of the same phenomenon. Attestors drawn from one interest class fail together,
and the standard model for that is design diversity, where independently produced versions are found
to exhibit coincident failures far more often than independence would allow; the review of
Littlewood, Popov and Strigini is the standard reference for those models, and Knight and Leveson's
experiment is the origin of the empirical objection \cite{littlewood2001,knight1986}. Read against
\eqref{eq:doubling}, that literature is measuring $\bar\rho$ in the sense of Remark
\ref{rmk:design-effect} rather than bounding it, which is what makes the framework's reliance on an
estimable companion quantity rather than on $q$ itself a practical choice and not merely a
concession.

\section{The Incentive Structure of the Intent Layer}\label{sec:mechanism}

Proposition \ref{th:correlated} attributes consensus reliability to the parameter $q$ but
gives no method of estimating $q$. This section identifies the nature of that omission: the
principal sources of $q$ are \emph{collusion} and \emph{shared interest}, so it is not a
statistical parameter but a problem of \emph{incentive design}. Appendix \ref{app:protocol} turns
the resulting requirements into a procedure and states its own failure modes; it is a sketch of an
unbuilt mechanism, not a validated one, and is offered so that the claim that this intersection is
unoccupied can be assessed against something concrete.

\begin{principle}[The composition of $q$]\label{pr:q-decomposition}
The common error mode decomposes into three sources:
\begin{enumerate}[label=(\alph*)]
\item \term{Cognitive isomorphism}: the verifiers share training provenance, domain
      assumptions or habits of thought;
\item \term{Interest isomorphism}: the verifiers' payoff is tied to ``passing'' rather than
      to ``being right'';
\item \term{Information isomorphism}: the verifiers read the same material under review, so
      their errors are determined by the same defect.
\end{enumerate}
The three sources correspond to three different interventions: admitting heterogeneous
participants (for (a)), changing the payoff structure (for (b)), and introducing independent
information channels (for (c)). Acting on the first alone is the common practice and the
least effective. \S\ref{sec:authority} identifies a fourth co-occurrence of the same family,
\emph{authority isomorphism}, and explains why it is not a fourth term of this decomposition but a
separate failure axis.
\end{principle}

\begin{principle}[Proper scoring rules]\label{pr:scoring}
For the second source, the correct remedy is to tie a verifier's reward to the
\emph{accuracy} of the judgement rather than to its \emph{lean}. A proper scoring rule
\cite{brier1950} supplies a ready-made instrument: the reward function attains its unique
maximum at the true belief, thereby removing the incentive for both ``going along with the
majority'' and ``pleasing the submitter''. A concrete form is a prediction market on the
proposition ``does this specification faithfully express the intent'', settled on the
defects and rework data observed subsequently. The mechanism has two by-products:
\begin{enumerate}[label=(\alph*)]
\item It provides a means of reducing $q$ other than replacing people;
\item It yields a posterior estimate of $q$: it suffices to measure the historical
      correlation of verifier errors.
\end{enumerate}
\end{principle}

This section closes the gap left by \S\ref{sec:consensus}. The bound there is conditional,
and its condition (independence) is here reduced to a quantity that can be designed,
measured and paid for. Chapter \ref{ch:program} writes it as hypothesis $H_{11}$.

\section{Who May Attest: the Deontic Axis}\label{sec:authority}

Sections \ref{sec:consensus} and \ref{sec:mechanism} treat attestation as an epistemic problem: is
the specification faithful to the intent, and how reliable is the majority that says so. A second
question exists, and the framework's own assurance condition already contains it without having
connected it. The second conjunct of \eqref{eq:assurance} requires that \emph{the contract is
explicit and authorised}. Authorisation is not a matter of reliability: a specification can be
faithfully attested, by a large and genuinely independent panel, and still be one that nobody had
the standing to adopt.

\begin{principle}[Authorisation is a separate failure axis]\label{pr:authority}
The reliability of an attestation network concerns whether its verdict is \emph{likely correct}.
Authorisation concerns whether the verdict \emph{decides anything}. The two are independent.
\begin{enumerate}[label=(\alph*)]
\item A verdict can be reliable and unauthorised: a competent panel drawn entirely from one
      stakeholder class, agreeing on a requirement that another class had the standing to refuse.
\item A verdict can be authorised and unreliable: a single accountable owner making an uninformed
      judgement.
\end{enumerate}
Consequently no improvement in $q$, in $n_{\mathrm{eff}}$, or in the incentive structure can repair an
authorisation failure. The failure is not that everyone is wrong together, but that everyone together
is not the body entitled to decide.
\end{principle}

This is why authority belongs \emph{beside} the three sources of $q$ rather than inside them.
Principle \ref{pr:q-decomposition} decomposes the common error mode into cognitive, interest and
information isomorphism---three ways for attestors to be wrong in the same way. A fourth
co-occurrence is easy to name, and call \term{authority isomorphism}: all attestors drawn from a
single stakeholder class. It does raise $q$, because one class tends to share both premises and
interests. But lowering $q$ does not address it. A perfectly diverse panel of outsiders is exactly
the case in which the authorisation failure is cleanest: they may be right and still lack standing.

The distinction has a well-worn institutional form, worth one sentence and no more. Distributed
ledger systems separate the mechanism that establishes \emph{what is true} from the mechanism that
establishes \emph{who may decide what counts as true}; consensus and governance are different
protocols, and a system with excellent consensus and no governance is not thereby legitimate. The
analogy is offered as an analogy. What it points at is a gap in this paper rather than a result in it.

\begin{remark}[Two consequences for the design of the intent layer]\label{rmk:authority-design}
\begin{enumerate}[label=(\alph*)]
\item The attestation record of \S\ref{sec:consensus} should carry the \emph{stakeholder class} of
      each attestor beside its provenance. That costs one field, makes authority isomorphism
      observable, and lets an authorisation failure be distinguished after the fact from a
      reliability failure---the two have different remedies.
\item The choice of the authorised body is not derivable from the framework. It is an institutional
      fact about who bears the consequences of the decision, and is fixed by the same argument that
      assigns $\Lk{\infty}$ to human bearers in Chapter \ref{ch:levels}---an argument whose premise
      is the prevailing liability regime, as that chapter states explicitly.
\end{enumerate}
\end{remark}

\section{Specification Binding}\label{sec:spec-binding}

The fourth pillar of \eqref{eq:assurance} binds the verified artefact to the executed artefact.
Nothing so far binds the specification to the intent it was meant to express, and the asymmetry
matters, because the two have opposite behaviour over time. A binary is byte-identical or it is not:
its binding is static and decidable once and for all. A specification is faithful to an intent that
moves. Proposition \ref{prop:three-grounds}(i) is the reason---requirements are generative, and
preferences are activated by situations that did not exist when the specification was written---so
fidelity is not a property a specification has but a quantity that decays.

Table \ref{tab:evidence-package} already requires item 10, re-verification rules, of every artefact:
which changes trigger recomputation of the affected evidence. The specification side has no
counterpart, and condition (i) of Proposition \ref{prop:assurance} is stated as though it needed none. The
repair is to make the specification carry its own binding record.

\begin{definition}[Specification binding]\label{def:spec-binding}
A \term{specification binding} is a four-tuple
\[
\mathrm{SB}(\Phi) \;=\; \bigl(\text{specification version},\; \text{attestation record},\;
\text{validity horizon},\; \text{re-confirmation triggers}\bigr),
\]
where the attestation record is the consensus artefact of \S\ref{sec:consensus} that established
$\Phi$'s fidelity when it was written, the validity horizon is the interval over which that
attestation is to be treated as current, and the re-confirmation triggers are the events that
invalidate it before the horizon expires.
\end{definition}

The triggers are the interesting part, and they are not arbitrary. Proposition \ref{prop:three-grounds}
names three sources of drift, and each suggests its own trigger.
\begin{enumerate}[label=(\alph*)]
\item \emph{Generative requirements}: a new class of user, scenario or deployment environment appears,
      which is precisely the situation that activates a preference nobody had expressed.
\item \emph{Lossy compression}: the domain model, ontology or controlled language in which $\Phi$ was
      written changes, so that the same text denotes something else.
\item \emph{Validation}: an observed incident whose resolution is not derivable from $\Phi$, which is
      evidence that $\Phi$ under-specifies the intent rather than that the implementation erred.
\end{enumerate}
Trigger (c) is worth singling out because it is available for free: an incident that required a human
judgement call is, by construction, an incident the specification did not settle.

Condition (i) of Proposition \ref{prop:assurance} should therefore be read with a time index. Writing
$\varepsilon_{\text{cons}}(\Delta t)$ for the consensus error after an interval $\Delta t$ since
attestation,
\begin{equation}\label{eq:spec-decay}
\Pr\bigl[\Phi \text{ is faithful to } I \text{ at time } t\bigr] \;\ge\;
1 - \varepsilon_{\text{cons}}(\Delta t),
\qquad \Delta t = t - t_{\text{attest}},
\end{equation}
which reduces to the original condition at $\Delta t = 0$. This paper does not give the form of
$\varepsilon_{\text{cons}}(\cdot)$: no decay model is proposed, and the quantity is left as a
measurable object rather than a modelled one. What the section establishes is the structural
symmetry---a specification needs a re-confirmation rule exactly as a binary needs a re-verification
rule---and that the symmetry is forced by Proposition \ref{prop:three-grounds}(i) rather than added
for completeness.

\begin{remark}[The specification-side dual of the artefact-binding gap]\label{rmk:spec-h9}
Hypothesis $H_9$ concerns the measurable risk of the artefact-binding gap. Its specification-side
dual is the risk of executing against a stale-but-attested specification: the attestation record is
still inside its horizon, the binary still matches its digest, and both are wrong together because
the intent moved. The two gaps are independent, and a deployment can close either while leaving the
other open.
\end{remark}

\section{The Consolidated Assurance Condition}\label{sec:assurance-final}

Combining the parts of this chapter yields a complete, conditional statement of assurance.

\begin{definition}[Two-layer trust structure]\label{def:two-layer}
The chain from intent to execution is a composite of a two-layer trust structure:
\begin{equation}\label{eq:twolayer}
\underbrace{I \xrightarrow{\;\text{consensus}\;} \Phi}_{\text{non-formalisable layer}}
\;\longrightarrow\;
\underbrace{\Phi \xrightarrow{\;\text{mechanical}\;} P \xrightarrow{\;\text{mechanical}\;} B
\xrightarrow{\;\text{mechanical}+\text{monitoring}\;} E}_{\text{formalisable layer}}
\end{equation}
In the formalisable layer, trust is guaranteed by mechanical verifiers; in the
non-formalisable layer, trust is guaranteed by adversarial consensus; and the boundary
between the two layers, the record of the consensus, and the traceability of the whole chain
are guaranteed mechanically by cryptographically anchored infrastructure.
\end{definition}

\begin{proposition}[Consolidated assurance: sufficient conditions]\label{prop:assurance}
Let $I$ be an intent, $\Phi$ an attested specification, $B$ a candidate implementation
(binary) and $E$ a runtime environment, with $A$ the declared environmental assumptions.
If the following conditions hold simultaneously, then ``property $\varphi$ holds within the
declared scope $\Sigma$'' holds:
\begin{align}
&\text{(i)}\quad \Pr\bigl[\Phi \text{ is faithful to } I \text{ at time } t\bigr] \;\ge\;
1 - \varepsilon_{\text{cons}}(\Delta t);\\
&\text{(ii)}\quad \Tr(B, E) \;\subseteq\; \Allow(\Phi);\\
&\text{(iii)}\quad E \models A;\\
&\text{(iv)}\quad \mathrm{digest}(B_{\text{deployed}}) = \mathrm{digest}(B_{\text{verified}});\\
&\text{(v)}\quad \varphi \text{ is preserved under the interpretation relation } R \text{ in use};\\
&\text{(vi)}\quad \Phi \text{ is authorised by a body with standing to adopt it,}
\end{align}
where $\varepsilon_{\text{cons}}(\cdot)$ is the bound of \eqref{eq:chernoff}, governed by $q$ and
$\bar\rho$, evaluated at $\Delta t = t - t_{\text{attest}}$.
\end{proposition}

The statement is deliberately one-directional. It was written elsewhere as an equivalence, and
that was wrong, for a reason worth stating because it recurs wherever a framework of this shape is
proposed. The six conditions are \emph{evidence} conditions: they say what must be in hand. They are
not truth conditions, and truth does not wait for them. A property can hold although no evidence for
it has been assembled, so the converse direction fails on the elementary ground that a system may be
correct and undocumented. What the proposition asserts is the useful direction---that holding all six
is enough---and what it deliberately does not assert is that holding fewer is impossible.

Each conjunct carries its own caveat, and none of them is mechanical in the same sense as the
others.
\begin{enumerate}[label=(\alph*)]
\item Condition (i) is a probability with no natural definition until faithfulness is defined, an
      adjudicator identified, and a class of errors specified. Until those exist, $\varepsilon$ is a
      parameter of a model rather than a measured frequency. Its reliability is governed by $q$ and
      by the correlation $\bar\rho$ of \S\ref{sec:consensus} and \S\ref{sec:mechanism}, and both are
      conditional-model quantities.
\item Condition (ii) is a semantic property, not a routine check. $\Tr(B,E) \subseteq \Allow(\Phi)$
      is decidable outright only for restricted languages, finite models or programs carrying
      checkable certificates; for general programs in general environments it is a proof obligation,
      discharged---when it is discharged at all---by a combination of typing, model checking,
      contract verification and testing of graded strength.
\item Conditions (iii) and (iv) are the mechanical pair, and their strength is limited by the trusted
      computing base and by the model of the environment, not by anything in this chapter.
\item Condition (v) is where the relational premise of Remark \ref{rmk:relational-compose} is spent.
\item Condition (vi) is deontic and is not a matter of reliability at all (\S\ref{sec:authority}).
\end{enumerate}
Three of the six conjuncts of Definition \ref{def:assurance} are therefore mechanical, one is
probabilistic, one is deontic, and one---the requirement that the contract be explicit---is what the
rest of the chapter is for.

One further distinction the consolidated statement must carry is the classical one between safety and
liveness. Condition (ii) constrains finite traces: nothing forbidden happens. Progress properties---a
request is eventually answered, a lock is eventually released, a queue eventually drains---are
liveness properties, they are not instances of ``no finite prefix violates $\Phi$'', and they are not
obtained by the same argument. Where a system's contract includes liveness, the obligations belong at
the levels and to the devices that can express them: fairness assumptions and well-foundedness
arguments for the protocol model, and monitoring with a bounded horizon for the deployed system. A
consolidated statement that quietly absorbed liveness into trace inclusion would be wrong in a way
that is easy to miss, and this one does not.

Equation \eqref{eq:twolayer} and Proposition \ref{prop:assurance} together give the complete
answer to Problem \ref{prob:trust}:

\begin{principle}[Answer to Problem \ref{prob:trust}]\label{pr:answer}
``Humans no longer read the code'' is attainable; ``humans need no longer answer for the
expression of the goal'' is not. The workable division of labour is: humans establish and
bear responsibility for the goal and the constraints; generators perform large-scale
generation, search and optimisation; mechanised verification and the execution boundary
judge and constrain implementations; and the system always carries traceable assurance
evidence.
\end{principle}

\begin{remark}[Summary of conclusions on the central questions Q1--Q8]\label{rmk:q1q8}
\begin{longtable}{L{5.4cm}L{9.0cm}}
\toprule
\textbf{Question} & \textbf{This paper's judgement} \\
\midrule
\endhead
Q1 Can it be proved that the intent and the specification agree? &
  Not by general automatic methods. Controlled languages, domain ontologies, examples and
  human attestation can narrow the problem, but cannot cover implicit requirements that were
  never expressed (Proposition \ref{prop:three-grounds}) \\
Q2 How can a specification be proved sufficient? &
  Consistency, satisfiability, realisability, vacuity and formal coverage metrics can all be
  checked; none of them proves that the specification contains every undeclared prohibition \\
Q3 How do local guarantees yield an end-to-end guarantee? &
  It requires that relations compose, that environmental assumptions close and that
  observation maps are compatible (Proposition \ref{prop:soundness}); liveness and real-time
  properties usually require additional conditions \\
Q4 Can functional and resource guarantees be unified? &
  Several constraints can be expressed in a joint behaviour--cost semantics, but time,
  memory, energy and information flow need different analysis models; joint expression is
  not unified solving \\
Q5 How can trust in the generative model be reduced? &
  The generator proposes candidates and an independent checker verifies them. This reduces
  dependence on the generator's correctness but does not remove the requirement to trust the
  logical kernel, the theorem statements, the external assumptions and the correspondence of
  artefacts \\
Q6 Can a given specification always be synthesised effectively? &
  No. Finite grammars, restricted logics and specific templates make some problems solvable;
  for arbitrary program semantic properties there are boundaries of undecidability and
  search-space explosion \\
Q7 Can actual latency be strictly guaranteed? &
  A usable worst-case execution time bound is obtainable under specific assumptions about
  hardware, inputs, scheduling and the cost model; benchmarks and asymptotic complexity are
  not by themselves sufficient to prove a bound at the nanosecond scale \\
Q8 Is a unified theory needed? &
  The better-founded choice is to establish shared proof interfaces and evidence formats
  while retaining the different semantics, assumptions and analysis tools of each domain. A
  unified formalism may be one goal, but it is not a necessary precondition \\
\bottomrule
\end{longtable}
\end{remark}

The empirical content of this chapter is directions 1--3 (Chapter \ref{ch:program}) and
hypotheses $H_7$, $H_8$ and $H_{11}$. Propositions \ref{prop:no-godel} and
\ref{prop:three-grounds} are arguments rather than theorems, and what can be challenged in
them is whether the three grounds are exhaustive. Theorem \ref{th:condorcet} and
Proposition \ref{th:correlated} are applications of standard probability theory, and their
conclusions depend on the independence assumption.
